\chapter*{参考资料与版本基线}
\addcontentsline{toc}{chapter}{参考资料与版本基线}

本册的机制解释以稳定概念为主；涉及 Linux 接口、调度策略、文件系统特性和持久化语义时，应回到带版本的主资料核对。下面列出本版技术校审使用的主要入口，核对日期为 2026 年 7 月 22 日。网页的“latest”内容会继续变化，做源码走读时仍须记录实际 checkout 的 tag 或 commit。

\begin{description}[leftmargin=2.4em,style=nextline]
\item[Linux 内核文档]
Linux 6.16 文档总入口：\url{https://www.kernel.org/doc/html/v6.16/}。

\item[公平调度]
EEVDF 调度器说明：\url{https://www.kernel.org/doc/html/latest/scheduler/sched-eevdf.html}；CFS 历史设计：\url{https://www.kernel.org/doc/html/latest/scheduler/sched-design-CFS.html}。

\item[用户指针访问]
\code{copy\_to\_user}/\code{copy\_from\_user} 返回值和上下文约束：\url{https://docs.kernel.org/6.11/kernel-hacking/hacking.html}。

\item[块设备粒度]
Linux block queue 的 logical block size、physical block size 与 minimum I/O size：\url{https://www.kernel.org/doc/html/v6.16/core-api/kernel-api.html}。

\item[持久化]
Linux \code{fsync(2)}，包括文件与父目录持久化边界：\url{https://man7.org/linux/man-pages/man2/fsync.2.html}。

\item[Btrfs]
在线碎片整理及其对 extent sharing 的影响：\url{https://btrfs.readthedocs.io/en/latest/Defragmentation.html}。

\item[F2FS]
F2FS 设计、checkpoint、后台 GC 与前台 GC：\url{https://www.kernel.org/doc/html/v6.16/filesystems/f2fs.html}。

\item[开放文件系统文档]
OpenZFS 文档：\url{https://openzfs.github.io/openzfs-docs/}；Microsoft 文件系统协议资料：\url{https://learn.microsoft.com/en-us/openspecs/windows_protocols/ms-fscc/efbfe127-73ad-4140-9967-ec6500e66d5e}；Apple File System Reference：\url{https://developer.apple.com/support/downloads/Apple-File-System-Reference.pdf}。
\end{description}

\begin{rubricbox}
引用的最低标准不是“网页存在”，而是每个易变化断言都能回答三个问题：针对哪个版本，适用哪些配置，原始资料具体支持哪一部分结论。性能数字还必须同时给出硬件、工作负载、队列深度、缓存与持久化设置。
\end{rubricbox}
